\subsection*{Tratamiento binario (cochera)}

\begin{table}[H]
\centering\small
\caption{Validación semisintética del estimador binario: sesgo, RMSE y cobertura (R6)}\label{tab:sim_b}
\resizebox{\textwidth}{!}{% Origen: notebook4_semisintetico.ipynb -> resultados/tablas/R6b_simulacion_binaria.csv (generado por informe/tablas_manuscrito.py)
\begin{tabular}{llrrrrr}
\toprule
Escenario & Estimador & $\theta$ verdadero & Sesgo & RMSE & Cobertura IC 95\% & Long. IC \\
\midrule
B1 & Hedónico OLS & 0.10 & 0.0002 & 0.0147 & 94 & 0.058 \\
B1 & DML PLR & 0.10 & -0.0040 & 0.0145 & 95 & 0.055 \\
B1 & DML-IRM (soporte común) & 0.10 & -0.0036 & 0.0198 & 96 & 0.081 \\
B2 & Hedónico OLS & 0.10 & 0.0412 & 0.0434 & 19 & 0.057 \\
B2 & DML PLR & 0.10 & 0.0392 & 0.0418 & 20 & 0.053 \\
B2 & DML-IRM (soporte común) & 0.10 & 0.0421 & 0.0462 & 45 & 0.078 \\
\bottomrule
\end{tabular}}
\origen{\texttt{notebook4\_semisintetico.ipynb} $\rightarrow$ \texttt{R6b\_simulacion\_binaria.csv}. 100 réplicas por escenario; N = 2\,051; efecto verdadero 0.10.}
\end{table}

\begin{figure}[H]
\centering
\includegraphics[width=0.85\textwidth]{F6b_semisintetico_binario.png}
\caption{Error del estimador binario en los escenarios B1 y B2 (F6).}\label{fig:sim_b}
\origen{\texttt{notebook4\_semisintetico.ipynb} $\rightarrow$ \texttt{F6b\_semisintetico\_binario.png}.}
\end{figure}

\paragraph{Sin confusor omitido (B1).} Los tres estimadores recuperan el efecto: el sesgo es de 0.000 a $-0.004$ y la cobertura del 94 al 96\,\% (tabla~\ref{tab:sim_b}). El IRM con soporte común tiene un IC algo más ancho (0.081 frente a 0.055 del PLR), el precio de no imponer una forma funcional. El hedónico lineal no se sesga en este caso, igual que en los datos reales, donde coincide con el IRM.

\paragraph{Con confusor moderado (B2).} Todos se sesgan en torno a $+0.04$, el 40\,\% del efecto verdadero, y la cobertura cae al 19--45\,\%. Ningún estimador corrige la confusión no observada. Además, la cota de sensibilidad evaluada con la fuerza real de $U$ fue de 0.024 en promedio, por debajo del sesgo real (0.042). El intervalo ampliado con la cota contuvo el efecto verdadero en solo el 1\,\% de las réplicas. \textbf{La cota es optimista para el estimador binario}: con el recorte al soporte común, la fuerza de $U$ medida a través de $\nu^2$ subestima su efecto real sobre el tratamiento. Por eso los benchmarks de la sección R5 deben leerse como cotas inferiores de la fragilidad.

\subsection*{Tratamiento continuo (baños)}

\begin{table}[H]
\centering\scriptsize
\caption{Validación semisintética del estimador de baños: sesgo, RMSE y cobertura del IC 95\,\% (\%)}\label{tab:sim}
\resizebox{\textwidth}{!}{% Origen: notebook4_semisintetico.ipynb -> resultados/tablas/R6_simulacion.csv (generado por informe/tablas_manuscrito.py)
\begin{tabular}{llrrrrrrr}
\toprule
Escenario & Estimador & $\theta$ verdadero & Sesgo & RMSE & Cob. iid & Cob. edificio & Cob. celda & Long. IC \\
\midrule
E1 & OLS sin ajustar & 0.0500 & 0.2867 & 0.2869 & -- & 0 & -- & 0.054 \\
E1 & Hedónico OLS (S1) & 0.0500 & 0.0251 & 0.0279 & -- & 50 & -- & 0.049 \\
E1 & DML PLR (S1) & 0.0500 & -0.0013 & 0.0125 & 95 & 95 & 94 & 0.047 \\
E1 & DML + precio/m$^2$ (inadmisible) & 0.0500 & -0.0500 & 0.0500 & 0 & 0 & -- & 0.006 \\
E2 & OLS sin ajustar & 0.0500 & 0.2917 & 0.2919 & -- & 0 & -- & 0.054 \\
E2 & Hedónico OLS (S1) & 0.0500 & 0.0351 & 0.0375 & -- & 24 & -- & 0.049 \\
E2 & DML PLR (S1) & 0.0500 & 0.0094 & 0.0162 & 83 & 84 & 85 & 0.047 \\
E3 & OLS sin ajustar & 0.0500 & 0.3036 & 0.3038 & -- & 0 & -- & 0.053 \\
E3 & Hedónico OLS (S1) & 0.0500 & 0.0580 & 0.0595 & -- & 1 & -- & 0.049 \\
E3 & DML PLR (S1) & 0.0500 & 0.0333 & 0.0359 & 22 & 22 & 23 & 0.047 \\
E4 & OLS sin ajustar & 0.0500 & 0.3230 & 0.3232 & -- & 0 & -- & 0.052 \\
E4 & Hedónico OLS (S1) & 0.0500 & 0.0956 & 0.0965 & -- & 0 & -- & 0.049 \\
E4 & DML PLR (S1) & 0.0500 & 0.0730 & 0.0743 & 0 & 0 & 0 & 0.047 \\
E5 & OLS sin ajustar & 0.0503 & 0.2822 & 0.2824 & -- & 0 & -- & 0.052 \\
E5 & Hedónico OLS (S1) & 0.0503 & 0.0247 & 0.0276 & -- & 51 & -- & 0.050 \\
E5 & DML PLR (S1) & 0.0503 & -0.0007 & 0.0125 & 96 & 95 & 95 & 0.047 \\
E6 & OLS sin ajustar & 0.0500 & 0.2611 & 0.2618 & -- & 0 & -- & 0.064 \\
E6 & Hedónico OLS (S1) & 0.0500 & 0.0214 & 0.0276 & -- & 64 & -- & 0.053 \\
E6 & DML PLR (S1) & 0.0500 & -0.0026 & 0.0181 & 81 & 88 & 94 & 0.056 \\
\bottomrule
\end{tabular}}
\origen{\texttt{notebook4\_semisintetico.ipynb} $\rightarrow$ \texttt{R6\_simulacion.csv}. 100 réplicas por escenario; N = 2\,220.}
\end{table}

\begin{figure}[H]
\centering
\includegraphics[width=\textwidth]{F6_semisintetico.png}
\caption{Validación semisintética de los baños. Izquierda: error por escenario; derecha: cobertura del IC del DML.}\label{fig:sim}
\origen{\texttt{notebook4\_semisintetico.ipynb} $\rightarrow$ \texttt{F6\_semisintetico.png}.}
\end{figure}

La tabla~\ref{tab:sim} confirma los resultados del entregable anterior sobre la muestra ROOFTOP:
\begin{itemize}[leftmargin=*]
\item \textbf{Sin confusión omitida (E1):} el DML es insesgado ($-0.001$) con una cobertura del 93\,\%. El hedónico lineal sobreestima en 0.022 y cubre el 50\,\%, porque no absorbe la confusión no lineal, lo que explica su diferencia con el DML en los datos reales. El control inadmisible lleva el efecto a cero.
\item \textbf{Con confusión omitida (E2--E4):} el DML se sesga (0.007, 0.025 y 0.054). En este caso la \textbf{cota de sensibilidad sí es fiel}: con la fuerza real de $U$ reproduce el sesgo (0.007, 0.023 y 0.050) y el intervalo ampliado cubre el 96, 98 y 95\,\%.
\item \textbf{Efecto heterogéneo (E5):} los GATE son insesgados (cobertura del 91--94\,\%) y la prueba de Wald detecta las diferencias el 49\,\% de las veces.
\item \textbf{Dependencia espacial (E6):} la cobertura iid es del 91\,\%, la agrupada por edificio del 91\,\% y la agrupada por bloque del 93\,\%.
\end{itemize}
